
The findings in Section \ref{sec:results} point to increases in wages, firm size, entry, and exit rates. These findings are consistent with increased business dynamism and concentration caused by the arrival of immigrants. I can use the model developed in Section \ref{sec:model} to explain the forces behind these changes as well as infer the type of these arriving immigrants, which are consistent with the empirical results. Firstly, one should note that the model is a partial equilibrium model in a static setting. Wages were taken exogenously, and remain an important channel through which many of these effects are driven. However, I will use economic intuition to argue that wages can be tied directly to managerial talent. A simple extension of the model can pin down the equilibrium wage given stationary distributions of incumbents and managerial talents. For now, I focus on Equation \eqref{eq:wage}, which shows the relation between wages and the managerial talent cut-off for new entrants. One can see that there is a positive relationship between wages and managerial talent. If one observes an increase in wages, this must be accompanied by an increase in said talent cut-off. In other words, the level of managerial talent one must have for it better to become an entrepreneur rather than a worker in the economy has increased. These effects are consistent with the story of the new immigrants being more talented than the existing population class. I can now paint a complete picture of the effects presented in Table \ref{tab:main}. Column (1) shows an increase in wages, which is interrelated to increased managerial talent in the economy. Column (2) shows increased firm size, which is a direct consequence of more talented managers demanding more labor on average (See Equation \eqref{eq:labordemand}). While our model makes no predictions about transition dynamics, if one expects that firm size and average managerial talent of managers have increased, then there must be some churn in the economy wherein older, less-talented incumbents exit. Simultaneously, one must expect these newly arriving higher talent migrants to set up new firms, which is reflected in the positive coefficient in column (3).

Table \ref{tab:bdsage} provides a clearer explanation of the story by allowing us to look at new firms and older incumbents. As discussed before, Panel A shows that this increased exit rate is among medium and old firms. Our empirical analysis is conducted in 5-year periods, while medium and old firms are 6+ years old. I argue that this reflects the fact that increased competition from more talented arriving incumbents leads to the exit of marginal incumbents. These were incumbent firms that were just above the talent cutoff to operate. When arriving immigrants push the wages up, these managers are no longer profitable and exit to become workers. To round things up, Panel B is consistent with the story of more talented arriving migrants. The firm size has increased across firms of all ages, i.e., the average talent of a manager has increased across the firms of all ages. To summarize, my empirical findings are consistent with the idea that arriving migrants are more talented than the existing stock of population in a county. These immigrants, in turn, increase the equilibrium talent cut-off for who becomes a manager, marginal incumbents exit, and firm size increases across the board. Importantly, all these effects are driven through the wage channel.

The findings of this study should be interpreted keeping in mind some of its limitations and simplifying assumptions. Firstly, as noted before, I reconcile the empirical effects of managerial talent through the wage channel without explicitly deriving this channel. However, it is easy to see the economic intuition behind it. Secondly, I do not account for any productivity growth associated with migration. Recent work has highlighted the importance of these effects, and shown that immigration is strongly connected to productivity growth in the US \citep{EffectofImmigrationProductivity, ImmigrationInnovationGrowth}. Finally, I do not account for any correlation or signal between capital endowment and managerial talent. It is highly plausible that these two are positively correlated. Work in the future can try to overcome some of these limitations to get a more robust understanding of entrepreneurship in the United States.  